\section{Causal Credibility Assessment}
\label{sec:pinning-credibility}

Chapter~\ref{chp:background} argued that empirical software engineering research advances most effectively when causal claims are accompanied by an explicit assessment of their credibility.
This section instantiates the four-step framework for the present study: derive a causal theory from existing literature, define the estimand and the assumptions under which it admits a causal interpretation, honestly acknowledge the limitations that bear on those assumptions, and navigate the alternative causal structures a careful reader would consider.
The goal is not to claim that the study's causal interpretation is beyond doubt, but to make the load-bearing assumptions transparent so that the strength of the evidence can be calibrated against them.

\subsection{Causal Theory}
\label{sec:pinning-credibility-theory}

Figure~\ref{fig:dag-pinning} shows the directed acyclic graph (DAG) governing RQ1.
\TODO{Draw \texttt{figs-pinning/dag-pinning.pdf} from the node inventory below; nodes inside the dashed box are absorbed by the design's project and time fixed effects, the node $G$ enters as an explicit covariate with an interaction, and nodes outside the box are residual threats.}

\begin{figure}[t]
    \centering
    \TODO{Insert dag-pinning.pdf here once drawn}
    \caption{Causal DAG for the effect of pinning direct dependencies on the five outcome metrics in RQ1.
    Nodes within the dashed region are time-invariant project characteristics absorbed by repository fixed effects; the time node is absorbed by month fixed effects; $G$ enters the regression as an explicit covariate with an interaction term.}
    \label{fig:dag-pinning}
\end{figure}

The treatment is the project-level versioning strategy for direct dependencies, contrasting all-pinning against the project's observed (predominantly floating) configuration.
The mechanism through which pinning affects the five outcome metrics is npm's deduplication algorithm: pinning fixes version constraints, which in turn changes the set of versions npm resolves into the transitive dependency graph, which in turn determines the value of each outcome metric.
The interaction between pinning and dependency-graph size $G$ is structural rather than incidental: the deduplication mechanism produces qualitatively different behavior at large $G$ because the probability of conflicting version constraints increases with graph complexity~\cite{DBLP:conf/icse/PinckneyCGBCG23}.

The DAG also names the confounders that, in an unstructured observational comparison of pinning and floating projects, would bias the estimated effect.
Project type (package versus application) shapes both the versioning strategy choice~\cite{DBLP:journals/tse/JafariCAST22} and the structural topology of the dependency graph~\cite{DBLP:conf/kbse/LatendresseMCS22, DBLP:journals/ese/DecanMG19}; library maintainers have strong incentives to use floating version ranges to avoid forcing duplicate installations on downstream consumers, while application maintainers have stronger incentives to pin for reproducibility \TODO{cite Wermke et al. 2025 lockfile design space, Renovate documentation}.
Project maturity --- expressed through age, release history, and contributor count --- predicts versioning strategy and outcome levels independently \TODO{cite Javan et al. TOSEM 2023, Zerouali et al. 2018}.
Calendar time matters because the npm ecosystem itself shifts over the observation window: the 2020 turn toward pinning, the rapid growth of the GitHub Advisory database~\cite{DBLP:journals/ese/ChinthanetKMIIM21}, and the changing prevalence of automated dependency-update tools all introduce time-varying common shocks.
Latent developer security-consciousness affects both the strategy choice and the speed of vulnerability remediation~\cite{DBLP:conf/ccs/PashchenkoVM20}.
Number of downstream dependents shapes maintainer conservatism through reputational pressure~\cite{DBLP:conf/uss/ZimmermannSTP19}.

The two main residual threats that the design does not fully absorb are time-varying tool adoption and upstream release frequency.
Adoption of Dependabot or Renovate within the study window changes the effective treatment a project receives --- a project floating in \Code{package.json} but using Renovate in pin mode is effectively pinned in deployment --- and independently reduces vulnerability lag \TODO{cite Dependabot impact study}.
Upstream release frequency affects both the maintenance burden of floating (and thus the choice to switch to pinning) and the rate at which dependencies become outdated.
Both are time-varying within projects and are not fully absorbed by project fixed effects.

The counterfactual simulation that underpins RQ1 fundamentally changes the identification problem relative to a standard observational comparison.
Because both potential outcomes $Y_{it}(0)$ and $Y_{it}(1)$ are constructed for every project-time unit from the same \Code{package.json}, the simulation severs every back-door path that runs from a confounder into the treatment assignment.
The remaining identification threats are not about which projects pin; they are about whether the simulated counterfactual outcomes have the same expected values as the actual potential outcomes would, and whether residual time-varying factors that the design does not absorb operate on the outcome side of the DAG.
The next subsection makes these threats precise.

\subsection{Causal Estimand and Identifying Assumptions}
\label{sec:pinning-credibility-estimand}

For each outcome metric $M$, let $Y_{it}(d)$ denote the value of $M$ for project $i$ at time $t$ under versioning strategy $d$, where $d=1$ denotes all-pinning and $d=0$ denotes the observed (predominantly floating) strategy.
The study targets two causal estimands.
The Average Treatment Effect, $\tau_{\text{ATE}} = E[Y_{it}(1) - Y_{it}(0)]$, captures the population-average effect of switching from the observed strategy to all-pinning.
The Conditional Average Treatment Effect by graph size, $\tau(G) = E[Y_{it}(1) - Y_{it}(0) \mid \text{size}(G_{it}) = G]$, captures how that effect varies with the size of the project's dependency graph.
The panel regression estimates the ATE through the main coefficient on \emph{pinning} and traces the CATE through the \emph{pinning} $\times \ln(\text{size}(G))$ interaction.
Because the simulation produces both potential outcomes for every project-time unit from the same \Code{package.json}, the design is structurally a within-project paired experiment, and the project and time fixed effects absorb time-invariant project heterogeneity and common time shocks.

Four identifying assumptions support the causal interpretation of the regression coefficients.

\paragraph{A1. Simulation fidelity.}
The simulated counterfactual outcomes have the same expected value as the actual potential outcomes would have had under the corresponding versioning strategy.
This is the primary identifying assumption: unlike in a standard observational study, selection-into-treatment confounding is not the principal threat --- measurement validity of the simulated counterfactual is.
A1 \emph{requires empirical argument}.
Violations arise if npm's resolver algorithm changed between the historical timestamp and the simulation date, or if local environment configurations differ between historical and simulated environments.

\paragraph{A2. No interference (SUTVA, Part 1).}
A project's outcomes under its assigned versioning strategy do not depend on the strategies adopted by other projects.
A2 is \emph{satisfied by construction}: dependency resolution is deterministic given the \Code{package.json} and a snapshot of the registry, so project $i$'s simulated treatment has no causal pathway to project $j$'s simulated outcome --- the registry state is fixed by the time-travel mechanism, and resolution does not consult other projects' configurations.
RQ2 deliberately violates A2 by design --- ecosystem-level interference is the entire point --- which is why RQ2 uses a network propagation simulation rather than a panel regression.

\paragraph{A3. Consistent treatment value (SUTVA, Part 2).}
There exists a single well-defined version of the treatment ``all-pinning.''
A3 is \emph{satisfied by construction}: the treatment is imposed by a deterministic transformation of the \Code{package.json} --- removing every caret and tilde from direct production dependencies and resolving to the minimum specified version --- yielding a single, operationally well-defined version of all-pinning.
The fact that real-world pinning is a spectrum bears on which estimand is policy-relevant, not on whether the simulated estimand is causally identified; that concern belongs to external validity and is addressed in the existing Limitations subsection (Section~\ref{sec:method-rq1}).

\paragraph{A4. No residual time-varying confounding.}
After conditioning on project fixed effects, time fixed effects, and $\ln(\text{size}(G_{it}))$, no time-varying confounder remains that affects both the effective treatment and the outcomes.
A4 \emph{requires empirical argument}.
The main candidate violator is automated dependency-update tool adoption (Dependabot, Renovate), which can change within the study window and affects the effective treatment definition (via Renovate's \Code{rangeStrategy}) and the outcomes for vulnerabilities, outdated dependencies, and automatic updates simultaneously.

The credibility of the causal claim therefore rests on A1 and A4 --- the two assumptions that the design's construction does not guarantee.
RQ2 is identified through a separate network propagation simulation rather than through the potential-outcomes framework above; its assumptions and their direction of bias are addressed in the existing limitations subsection of the RQ2 method section (Section~\ref{sec:method-rq3}).

\subsection{Limitations and Known Caveats}
\label{sec:pinning-credibility-limitations}

The simulation design eliminates the standard observational confounders identified in the DAG --- selection into the pinning strategy by project type, maturity, security-consciousness, CI adoption, and downstream dependents --- and the panel fixed effects absorb time-invariant project heterogeneity and common time shocks.
Two threats to the internal causal validity of the RQ1 estimate remain.

\paragraph{Differential simulation failures.}
Resolution failures occurred for 12--16\% of project-time observations and are not random.
Pinning introduces additional version conflicts, so success rates differ between the treatment and control conditions: projects whose dependency constraints already conflict are more likely to fail under pinning than under the observed (predominantly floating) configuration.
The projects with the most severe latent conflicts --- precisely those for which the bloat cost of pinning would be highest --- are disproportionately excluded from the analysis sample, which threatens A1.
The reported pinning effect on \Code{n\_bloated} is therefore a lower bound on the true cost of pinning, and the magnitude of the underestimation increases with $G$.

\paragraph{Time-varying confounding by automated dependency-update tool adoption.}
Adoption of Dependabot or Renovate is not uniformly time-invariant within the 12-month panel; some projects adopt these tools mid-window.
This adoption confounds the relationship between the declared versioning strategy and the outcomes for vulnerabilities, outdated dependencies, and automatic updates: projects using Renovate in pin mode have an effective treatment that differs from their \Code{package.json} alone, and these tools independently reduce vulnerability lag.
The threat is to A4.
The direction of bias depends on whether tool adoption correlates with $G$ within projects; if security-conscious projects both adopt these tools and have larger dependency graphs, the $G$ interaction term partially absorbs the confounding.

The remaining limitations of the study --- the all-or-nothing definition of the simulated treatment, the absence of a modeled developer behavioral response, the uniform random attack assumption underlying the interpretation of \Code{n\_floating} as an attack surface, the npm-specific scope of the deduplication mechanism, and the simplifying assumptions in the RQ2 ecosystem-level network simulation --- are concerns about external validity, measurement validity, treatment definition, or ecosystem-specific scope rather than threats to the internal identification of the RQ1 causal estimate.
They are addressed in the existing ``Limitations and Threats to Validity'' paragraphs in Section~\ref{sec:method-rq1} (RQ1) and Section~\ref{sec:method-rq3} (RQ2).

\subsection{Alternative Explanations}
\label{sec:pinning-credibility-alternatives}

Three alternative causal structures could in principle explain the central RQ1 findings without the proposed mechanism; this subsection considers each and explains how the design rules them out, partially addresses them, or leaves them open.

\paragraph{Alternative 1: The crossover finding is a simulation artefact.}
The counter-intuitive result that pinning increases \Code{n\_floating} for large dependency graphs (Figure~\ref{fig:n-floating}) could be a simulation artefact if the \Code{-{}-before} resolver produces systematically different deduplication behavior for pinned versus floating \Code{package.json} files at high graph complexity.
The design is defensible against this alternative: the mechanism --- deduplication producing additional transitive version copies when pinning forces version conflicts --- is structurally documented in npm's resolver implementation~\cite{how-npm3-works} and is illustrated with a minimal worked example in Figure~\ref{fig:rq2-example}.
The crossover point at 498 nodes is consistent with the mechanism activating only once enough version conflicts accumulate in the transitive graph, a continuous rather than discontinuous process.
The finding replicates across four robustness checks --- npm packages only, GitHub repositories only, $t_0$ only, and development dependencies included --- each representing a distinct subpopulation.
A simulation artefact would need to produce the same spurious crossover point in all four populations, which is implausible.

\paragraph{Alternative 2: Time-varying confounding by automated tool adoption inflates the cost estimates.}
Projects that adopted Dependabot or Renovate during the study window have lower \Code{n\_vuln} and \Code{n\_outdated\_deps} regardless of their pinning strategy.
If tool adoption increased over the observation window and correlates with $G$, the estimated cost of pinning on these two outcomes could be confounded.
The design partially addresses this: project fixed effects absorb time-invariant tool adoption, leaving only the new adoptions within the 12-month window as a residual threat.
Tool adoption tends to be lumpy rather than continuously varying --- projects either have these tools or they do not for the duration of the panel --- which limits the magnitude of the residual confounding.
The direction, if present, would reduce the estimated costs of pinning, since automated tools compensate for the vulnerability lag pinning induces; the reported costs are therefore conservative lower bounds in this respect.

\paragraph{Alternative 3: Unobserved selection on project quality drives the within-project comparison.}
A reader could ask whether projects whose maintainers chose pinning are systematically different in unobserved ways --- more security-conscious, better resourced, more careful about updates --- and whether those differences, rather than the pinning intervention itself, drive the observed effects.
The counterfactual simulation eliminates this concern at the level of the RQ1 estimate: treatment is imposed by a deterministic transformation of the \Code{package.json}, not chosen by the project's maintainer.
Selection into the observed pinning strategy is irrelevant to the simulated potential outcomes because every project is observed under both treatment conditions.
This is the principal advantage of the counterfactual simulation design over an observational comparison and is what permits the relatively strong causal claims of RQ1 despite the absence of randomization.
